\chapter{Playing and monitoring}\label{ch:playing}

\section{Play commands}

\begin{center}
\begin{tabular}{@{}ll@{}}
\toprule
Key & Action \\ \midrule
\key{F5} & play the song from the start \\
\key{F7} & play from the cursor \\
\key{SHIFT+F7} & play from the bookmark \\
\key{F6} & play the current song line in a loop \\
\key{SHIFT+F6} & play the selected block in a loop \\
\key{ESC} & stop and silence \\
\key{F12} & follow mode: the screen follows the playing position \\
\key{CONTROL+F12} & toggle PAL / NTSC \\
\key{F11} & respect volume on/off \\
\bottomrule
\end{tabular}
\end{center}

The same commands are the buttons of the main toolbar and the entries of the \menu{Play} menu. With the option \emph{Reset the Atari sound
routines each time ESC is pressed} (on by default) \key{ESC} also starts the player routine again, which clears sounds that hang.
A song with the AUTOFILTER that is opened while the Patch16 driver is selected makes the program ask whether to play it with the
original RMT 1.28 driver (chapter~\ref{ch:drivers}).

\section{Channels}

The \menu{Channels} menu and the keys switch the channels on and off to hear a part alone: \key{F9} mutes or unmutes the active
channel, \key{SHIFT+F9} all the channels, \key{CONTROL+F9} plays the active channel alone (solo) and again unmutes everything,
\key{CONTROL+1}--\key{8} toggle a channel. The mouse does the same on the track headers (chapter~\ref{ch:window}).

\section{Watching the sound}

\begin{itemize}
  \item The \emph{play time counter} (\menu{View $\rightarrow$ Play Time Counter}) shows the time and the tempo while the song plays.
  \item The \emph{volume analyzer} (\menu{View $\rightarrow$ Volume Analyzer}) draws the volume of each channel above the track
        headers; figure~\ref{fig:window-analyzer} shows it.
  \item \menu{Pokey $\rightarrow$ Pokey Chip Registers} shows the registers of the POKEY with the pitch, volume and distortion
        of each channel: the real thing the Atari would do.
\end{itemize}

\shot{window-analyzer}{The volume analyzer and the registers of the POKEY while a song plays.}{0.98}

\section{Playing from the keyboard}

In prove mode (\key{CONTROL+SPACE}) the note keys only play. In edit mode \key{SHIFT} plus a note key plays without writing. A
MIDI keyboard does the same and can also record (chapter~\ref{ch:midi}).
